% =========================================================
% Pakete
% =========================================================

% --- Basis / Sprache / Encoding ---
\usepackage[ngerman]{babel}            % Sprachdefinition für alles
\usepackage[T1]{fontenc}
\usepackage[utf8]{inputenc}           % UTF-8 Encoding (nur nötig bei älteren Engines)
\usepackage[full]{textcomp}                 % Verhindert Warnungen bei bestimmten Sonderzeichen
\usepackage{csquotes}                 % nötig für biblatex und \textquote (Anführungszeichen)
\usepackage{etoolbox}                 % programming facilities (z.B. Symbolverzeichnis)

% --- Literatur / Quellen ---
\usepackage[backend=biber,style=ieee]{biblatex} % Bibliotheksmanagement

% --- Schrift / Typografie ---
\usepackage{lmodern}
\usepackage{opensans}                 % zusätzliche Schrift (wie in deinem neuen Block)
\usepackage[scaled]{helvet}           % Lädt Helvetica (phv)
\renewcommand{\familydefault}{\sfdefault}
\usepackage[italic]{mathastext}
\Mathastext

% --- Layout / Seitenränder / Kopf-Fuß ---
\usepackage[a4paper,tmargin=2cm,lmargin=3.2cm,rmargin=2cm,bmargin=3cm]{geometry} % Satzspiegel
\usepackage{fancyhdr}                 % Kopf-/Fußzeile
\usepackage{extramarks}
\usepackage{lastpage}                 % Ausgabe der letzten Seite

% --- Mathe / Einheiten ---
\usepackage{amsmath,amssymb}          % Mathe + Symbole (amssymb neu)
\usepackage{stix}                     % hochwertige Symbole & Schrift (neu)
\usepackage{siunitx}                  % Einheiten (neu)
\usepackage{xfrac}                    % Brüche im Text (z.B. \sfrac{1}{2})
\usepackage{icomma}                   % kein Platz nach Komma in mathmode
\sisetup{
    locale = DE,
    per-mode = symbol
}
\DeclareSIUnit{\pct}{\%}

% --- Grafiken / PDFs / SVG ---
\usepackage{graphicx}
\graphicspath{{Bilder/}}              % Bildpfad
\usepackage{subcaption}               % Subfigures (neu)
\usepackage{pdfpages}                 % PDFs als ganze Seiten einbinden
\usepackage{svg}                      % SVG-Bilder
\usepackage{afterpage}

% --- Zeichnungen ---
\usepackage{tikz}
\usepackage{calc}
\usetikzlibrary{calc}
\usetikzlibrary{arrows.meta,positioning} % neu (aus deinem neuen Block)
\newcommand{\cellnode}[1]{\tikz[remember picture,baseline]\node (#1) {};}

% --- Farben ---
\usepackage[table]{xcolor}            % Farben (u.a. Tabellen)

% RGB definitions (print-safe, pure RGB) - neu
\definecolor{conceptA}{RGB}{220,40,40}   % Red
\definecolor{conceptB}{RGB}{40,160,80}   % Green
\definecolor{conceptC}{RGB}{40,80,200}   % Blue

% --- Tabellen ---
\usepackage{tabularx}
\usepackage{longtable}               % neu
\usepackage{multirow}                % neu
\usepackage{tablefootnote}           % Fußnoten in Tabellen
\usepackage{pdflscape}

% --- Float-/Platzierung ---
\usepackage{float}                   % neu
\usepackage[section]{placeins}        % \FloatBarrier
\usepackage{array}

% --- Inhalts-/Abbildungs-/Tabellenverzeichnisse ---
\usepackage[nottoc,notbib]{tocbibind}

% --- Caption / Fußnotenformat ---
\usepackage{caption}
% --- Caption Optionen sicher setzen (muss VOR subcaption/caption kommen)
\PassOptionsToPackage{hang,font=small}{caption}

%\usepackage[hang,font=small]{caption} % Mehrzeilige Darstellung von Überschriften
\usepackage[hang]{footmisc}           % Mehrzeiliges Einrücken von Fußnoten
\setlength{\belowcaptionskip}{12pt}   % neu (aus deinem neuen Block)

% --- Positionierung (z.B. Titelseite) ---
\usepackage[absolute]{textpos}

% --- Sonstiges ---
\usepackage[normalem]{ulem}           % \sout, Unterstreichen etc.
\usepackage{comment}                 % neu (aus deinem neuen Block; oben war es auch schon mal)
%\usepackage{todonotes}

% -------------------------------------------------
% Code Listings
% -------------------------------------------------
\usepackage{listings}
\usepackage{xcolor}
\usepackage[T1]{fontenc}
\usepackage[utf8]{inputenc}
\usepackage{textcomp}

\lstset{
  inputencoding=utf8,
  extendedchars=true
}

% -------------------------------------------------
% VS Code–inspired colors (adapted for white background)
% -------------------------------------------------
\definecolor{vscKeyword}{RGB}{175,0,219}
\definecolor{vscString}{RGB}{163,21,21}
\definecolor{vscComment}{RGB}{0,128,0}
\definecolor{vscNumber}{RGB}{9,134,88}
\definecolor{vscFunction}{RGB}{43,145,175}
\definecolor{vscBg}{RGB}{245,245,245} % very light gray background

% -------------------------------------------------
% Python listing style
% -------------------------------------------------
\lstdefinestyle{python-vsc}{
    language=Python,
    basicstyle=\ttfamily\small,
    keywordstyle=\color{vscKeyword},
    stringstyle=\color{vscString},
    commentstyle=\color{vscComment},
    numberstyle=\tiny\color{gray},
    numbers=left,
    stepnumber=1,
    numbersep=8pt,
    backgroundcolor=\color{vscBg},
    frame=single,
    rulecolor=\color{gray!40},
    framesep=6pt,
    frameround=ffff,
    breaklines=true,
    showstringspaces=false,
    tabsize=4,
    captionpos=b,
    xleftmargin=2.5em,
    framexleftmargin=2.0em
}

% -------------------------------------------------
% TMCL / Trinamic listing style (pseudo language via C)
% -------------------------------------------------
\lstdefinestyle{tmcl-vsc}{
    language=C,
    basicstyle=\ttfamily\small,
    keywordstyle=\color{vscKeyword},
    stringstyle=\color{vscString},
    commentstyle=\color{vscComment},
    numberstyle=\tiny\color{gray},
    numbers=left,
    numbersep=8pt,
    backgroundcolor=\color{vscBg},
    frame=single,
    rulecolor=\color{gray!40},
    framesep=6pt,
    frameround=ffff,
    breaklines=true,
    showstringspaces=false,
    tabsize=4,
    captionpos=b,
    xleftmargin=2.5em,
    framexleftmargin=2.0em
}

% -------------------------------------------------
% Hyperlinks (should be loaded late)
% -------------------------------------------------
\usepackage[
  colorlinks=true,
  linkcolor=black,
  citecolor=black,
  urlcolor=black
]{hyperref}


% --- Hyperlinks (sollte weit hinten stehen) ---
\usepackage[colorlinks=true,linkcolor=black,citecolor=black,urlcolor=black]{hyperref}


% =========================================================
% Text- und Schriftgrößen / Umbruchsregeln
% =========================================================
\setlength{\parindent}{0em}           % kein Einrücken
\newcommand{\changefont}{\fontsize{11}{13}\selectfont}      % Kopfzeile
\newcommand{\changefontbig}{\fontsize{10}{12}\selectfont}  % Kopfzeile

\setcounter{biburllcpenalty}{9000}    % Umbruchverhalten URLs (Kleinbuchstaben)
\setcounter{biburlucpenalty}{9000}    % Umbruchverhalten URLs (Großbuchstaben)

\widowpenalty = 10000                 % Vermeidung von Satzende auf nächster Seite
\clubpenalty  = 10000                 % Vermeidung von Absatzende auf nächster Seite


% =========================================================
% Nummerierungen
% =========================================================
\numberwithin{equation}{section}
\numberwithin{table}{section}
\numberwithin{figure}{section}


% =========================================================
% Listenformat (academic style)
% =========================================================
\renewcommand{\labelenumii}{\arabic{enumi}.\arabic{enumii}}
\renewcommand{\labelenumiii}{\arabic{enumi}.\arabic{enumii}.\arabic{enumiii}}
\renewcommand{\labelenumiv}{\arabic{enumi}.\arabic{enumii}.\arabic{enumiii}.\arabic{enumiv}}
\renewcommand{\labelitemi}{-}


% =========================================================
% Kopfzeile: Section links, Seitenzahl rechts (article)
% =========================================================

\setlength{\headheight}{15pt}

\pagestyle{fancy}
\fancyhf{}

% Set running mark from \section (number + title)
\renewcommand{\sectionmark}[1]{\markright{\thesection\ #1}}

% Header content
\fancyhead[L]{\changefont \nouppercase{\rightmark}}
\fancyhead[R]{\changefont \thepage}

\renewcommand{\headrulewidth}{0.4pt}
\renewcommand{\footrulewidth}{0pt}

% Apply also on "plain" pages (some environments force plain)
\fancypagestyle{plain}{%
  \fancyhf{}%
  \fancyhead[L]{\changefont \nouppercase{\rightmark}}%
  \fancyhead[R]{\changefont \thepage}%
  \renewcommand{\headrulewidth}{0.4pt}%
}

% =========================================================
% autoref: Anpassungen (Equations + Kapitel-/Abschnittsnamen)
% =========================================================

% autoref für Equation: "Gleichung (x)"
\makeatletter
  \AtBeginDocument{%
    \let\plain@equationautorefname\equationautorefname
    \def\equationautorefname{\plain@equationautorefname\@autoref@insert@tagform}%
    \def\@autoref@insert@tagform~#1\null{~(#1\null)}%
  }
\makeatother

% Kapitelverweise / Abschnittsnamen
\addto\extrasngerman{\def\subsectionautorefname{Unterkapitel}}
\addto\extrasngerman{\def\subsubsectionautorefname{Unterkapitel}}
\addto\extrasngerman{\def\sectionautorefname{Kapitel}}

% (zusätzlicher Block aus deinem neuen File: subsection/subsubsection = Abschnitt)
% -> NICHT entfernt; aber beachte: das überschreibt ggf. die oberen "Unterkapitel"
\addto\extrasngerman{
  \def\subsectionautorefname{Abschnitt}
  \def\subsubsectionautorefname{Abschnitt}
}

% Listings
\renewcommand{\lstlistingname}{Code-Ausschnitt}
\def\lstnumberautorefname{Zeile}


% =========================================================
% Listings: Spyder-Theme Farben + Stil
% =========================================================

\definecolor{spyderBlue}{HTML}{0000FF}
\definecolor{spyderGreen}{HTML}{008000}
\definecolor{spyderGrey}{HTML}{999999}
\definecolor{spyderCyan}{HTML}{9a0aab}
\definecolor{spyderNumber}{HTML}{B8860B}
\definecolor{spyderself}{HTML}{b80b1f}

\lstdefinestyle{spyderstyle}{
    language=Python,
    backgroundcolor=\color{white},
    commentstyle=\color{spyderGrey},
    keywordstyle=\color{spyderCyan},
    stringstyle=\color{spyderGreen},
    numberstyle=\tiny\color{darkgray},
    basicstyle=\ttfamily\footnotesize\color{black},
    emphstyle=\color{spyderCyan},
    emph={[2]int, str, float, list, dict, set, tuple, print, range, open, len},
    emphstyle={[2]\color{spyderBlue}\bfseries},
    emph={[3]self},
    emphstyle={[3]\color{spyderself}\bfseries},
    escapeinside={(*@}{@*)},
    literate=
      *{0}{{{\color{spyderNumber}0}}}1
       {1}{{{\color{spyderNumber}1}}}1
       {2}{{{\color{spyderNumber}2}}}1
       {3}{{{\color{spyderNumber}3}}}1
       {4}{{{\color{spyderNumber}4}}}1
       {5}{{{\color{spyderNumber}5}}}1
       {6}{{{\color{spyderNumber}6}}}1
       {7}{{{\color{spyderNumber}7}}}1
       {8}{{{\color{spyderNumber}8}}}1
       {9}{{{\color{spyderNumber}9}}}1
       {.}{{{\color{spyderNumber}.}}}1
       {e-}{{{\color{spyderNumber}e-}}}2
       {e+}{{{\color{spyderNumber}e+}}}2
       {E-}{{{\color{spyderNumber}E-}}}2
       {E+}{{{\color{spyderNumber}E+}}}2,,
    breakatwhitespace=false,
    breaklines=true,
    captionpos=b,
    keepspaces=true,
    numbers=left,
    numbersep=5pt,
    showspaces=false,
    showstringspaces=false,
    showtabs=false,
    tabsize=4
}

\lstset{style=spyderstyle}


% =========================================================
% Anhangsverzeichnis (Haupt-TOC + Anhangs-TOC)
% =========================================================

\makeatletter
\newcommand*{\maintoc}{%
  \begingroup
    \@fileswfalse
    \renewcommand*{\appendixattoc}{\value{tocdepth}=-10000}
    \tableofcontents
  \endgroup
}

\newcommand*{\appendixtoc}{%
  \begingroup
    \edef\@alltocdepth{\the\value{tocdepth}}
    \setcounter{tocdepth}{-10000}
    \renewcommand*{\contentsname}{Anhangsverzeichnis}
    \renewcommand*{\appendixattoc}{\setcounter{tocdepth}{\@alltocdepth}}
    \tableofcontents
    \setcounter{tocdepth}{\@alltocdepth}
  \endgroup
}

\newcommand*{\appendixattoc}{}
\g@addto@macro\appendix{%
  \clearpage
  \addcontentsline{toc}{section}{\appendixname}
  \addtocontents{toc}{\protect\appendixattoc}
}
\makeatother


% =========================================================
% Kurzbefehle / eigene Commands
% =========================================================

% Römische Zahlen
\newcommand{\uproman}[1]{\uppercase\expandafter{\romannumeral#1}}
\newcommand{\lowroman}[1]{\romannumeral#1\relax}

% plus und minus in kleiner
\newcommand{\plus}{\raisebox{.4\height}{\scalebox{.6}{+}}}
\newcommand{\minus}{\raisebox{.4\height}{\scalebox{.8}{-}}}

% Kreis um z.B. Zahl
\newcommand\circlearound[1]{%
  \unitlength1ex
  \begin{picture}(2.5,2.5)%
    \put(0.75,0.75){\circle{2.5}}%
    \put(0.75,0.75){\makebox(0,0){#1}}%
  \end{picture}%
}


% =========================================================
% Kurzbefehle fürs Einfügen von Bildern und Svgs
% =========================================================

\newcommand*{\newfigure}[4][]{%
    \begin{figure}[ht!]%
        \centering%
        \includegraphics[width=#2\textwidth]{#3}%
        \caption[#1]{#4}%
        \label{fig:#3}%
        \ifx&#1&%
            \addcontentsline{lof}{figure}{\protect\numberline{\thefigure}#4}%
        \fi%
    \end{figure}%
    \FloatBarrier
}

\newcommand*{\newsvg}[4][]{%
    \begin{figure}[ht!]%
        \centering%
        \includesvg[width=#2\textwidth]{#3}%
        \caption[#1]{#4}%
        \label{fig:#3}%
        \ifx&#1&%
            \addcontentsline{lof}{figure}{\protect\numberline{\thefigure}#4}%
        \fi%
    \end{figure}%
    \FloatBarrier
}

\newcommand*{\newmultifig}[4]{%
\begin{figure}[ht!]
\centering
\begin{minipage}[t]{0.46\textwidth}
    \centering
    \includegraphics[width=\textwidth]{#1}
    \caption{#2}
    \label{fig:#1}
\end{minipage}
\hspace{0.03\textwidth}
\begin{minipage}[t]{0.46\textwidth}
    \centering
    \includegraphics[width=\textwidth]{#3}
    \caption{#4}
    \label{fig:#3}
\end{minipage}
\end{figure}
\FloatBarrier
}

\newcommand{\newparagraph}[1]{%
\paragraph{#1}\label{#1}\mbox{}\\[1ex]
}
